\chapter{Graphical Criteria for Identification}
\label{chap:graphical}

\section{Causal diagrams}
\label{sec:graph-dags}

A causal directed acyclic graph\index{directed acyclic graph} encodes a set of
structural assumptions as a picture. Vertices are variables, a directed edge
$U \to V$ asserts that $U$ may be a direct cause of $V$, and the absence of an
edge is the substantive claim: it asserts that no direct effect exists. A graph
with every possible edge present says nothing, which is a useful reminder that
the content of a diagram lies in what it omits.

\begin{definition}[Causal DAG]
\label{def:dag}
A causal DAG over variables $V = \{V_1, \dots, V_p\}$ is a directed acyclic
graph $G$ together with the assertion that each $V_j$ is generated as
$V_j = f_j\left( \mathrm{pa}_G(V_j), \epsilon_j \right)$ with mutually
independent disturbances $\epsilon_j$, where $\mathrm{pa}_G(V_j)$ denotes the
parents of $V_j$ in $G$.
\end{definition}

Figure~\ref{fig:confounding-dag} shows the structure assumed for the retraining
study of Section~\ref{sec:po-cost}. Prior earnings and regional labour demand
affect both enrolment and later earnings. Motivation is unmeasured and affects
both. A caseworker referral affects enrolment but, by assumption, not earnings
except through enrolment.

\begin{figure}[t]
\centering
\begin{tikzpicture}[
  >={Stealth[length=2.2mm]},
  node distance=17mm,
  obs/.style={draw,circle,minimum size=9mm,inner sep=0pt,font=\small},
  unobs/.style={draw,circle,dashed,minimum size=9mm,inner sep=0pt,font=\small},
  every edge/.style={draw,->,thick}
]
  \node[obs] (A) {$A$};
  \node[obs,right=26mm of A] (Y) {$Y$};
  \coordinate (mid) at ($(A)!0.5!(Y)$);
  \node[obs] (X) at ($(mid)+(0,15mm)$) {$X$};
  \node[unobs] (U) at ($(mid)-(0,15mm)$) {$U$};
  \node[obs,left=20mm of A] (Z) {$Z$};

  \draw (A) edge (Y);
  \draw (X) edge (A);
  \draw (X) edge (Y);
  \draw (U) edge (A);
  \draw (U) edge (Y);
  \draw (Z) edge (A);

  \node[font=\footnotesize,above=1mm of X] {prior earnings, regional demand};
  \node[font=\footnotesize,below=1mm of U] {motivation, unmeasured};
  \node[font=\footnotesize,below=1mm of Z,align=center] {caseworker\\ referral};
  \node[font=\footnotesize,below=1mm of A] {enrolment};
  \node[font=\footnotesize,below=1mm of Y] {earnings};
\end{tikzpicture}
\caption[Causal diagram for the retraining study]{Causal diagram for the
retraining study. Solid vertices are measured, the dashed vertex is not. The
covariate set $X$ blocks the backdoor path through prior earnings and regional
demand. It does not block the path through motivation, which is why
Chapter~\ref{chap:sensitivity} asks how strong $U$ would have to be to overturn
the conclusion, and why $Z$ is available as an instrument in
Section~\ref{sec:design-iv}.}
\label{fig:confounding-dag}
\end{figure}

\section{d-separation and the backdoor criterion}
\label{sec:graph-backdoor}

\begin{definition}[Blocked path]
\label{def:blocked}
\index{d-separation}
A path $\pi$ between $A$ and $Y$ is blocked by a set $S$ if it contains a chain
$V_{j-1} \to V_j \to V_{j+1}$ or a fork $V_{j-1} \leftarrow V_j \to V_{j+1}$
with $V_j \in S$, or a collider $V_{j-1} \to V_j \leftarrow V_{j+1}$ with
neither $V_j$ nor any descendant of $V_j$ in $S$.
\end{definition}

\begin{definition}[Backdoor criterion]
\label{def:backdoor}
\index{backdoor criterion}
A set $S$ satisfies the backdoor criterion relative to $(A, Y)$ in $G$ if no
element of $S$ is a descendant of $A$, and $S$ blocks every path between $A$ and
$Y$ that begins with an edge into $A$.
\end{definition}

\begin{theorem}[Backdoor adjustment]
\label{thm:backdoor}
If $S$ satisfies the backdoor criterion relative to $(A,Y)$ in a causal DAG $G$,
and positivity holds for $S$, then
\begin{equation}
  \E\left[ Y(a) \right] = \int \E\left[ Y \mid A = a, S = s \right] \, dF_S(s).
  \label{eq:backdoor-adjust}
\end{equation}
\end{theorem}

\begin{proof}
The structural model of Definition~\ref{def:dag} implies the local Markov
property\index{Markov property}, so the joint density factorises as
$f(v) = \prod_j f\left( v_j \mid \mathrm{pa}_G(v_j) \right)$. The intervention
$do(A = a)$ replaces the factor $f(a \mid \mathrm{pa}_G(A))$ by a point mass,
giving the truncated factorisation
$f_a(v) = \prod_{j \neq A} f\left( v_j \mid \mathrm{pa}_G(v_j) \right)$.
Marginalising $f_a$ over all variables other than $Y$ and $S$, and using that
$S$ contains no descendant of $A$ so that its marginal law is unaffected by the
truncation, yields $f_a(y) = \int f(y \mid a, s) \, dF_S(s)$. The backdoor
condition ensures $Y(a) \indep A \mid S$, which is what licenses replacing
$f(y \mid do(a), s)$ by $f(y \mid a, s)$. Taking expectations
gives~\eqref{eq:backdoor-adjust}.
\end{proof}

Theorem~\ref{thm:backdoor} and Theorem~\ref{thm:ate-identification} are the same
result in two languages. The graphical version has the practical advantage that
the adjustment set is derived from the assumed structure rather than asserted,
and the disadvantage that the structure must be asserted first.

\begin{example}[A collider that should not be adjusted for]
\label{ex:collider}
\index{collider bias}
Suppose the analyst adds an indicator of programme completion to the adjustment
set for Figure~\ref{fig:confounding-dag}. Completion is a descendant of
enrolment and is also affected by motivation, so it is a collider on the path
$A \to C \leftarrow U \to Y$. Conditioning on it opens that path and induces an
association between enrolment and unmeasured motivation within completion
strata. The resulting bias has no sign that can be signed in advance, and adding
more covariates does not reduce it.
\end{example}

Example~\ref{ex:collider} is the reason that ``adjust for everything measured''
is not a defensible default \citep{sorensen2013collider}. In the retraining data, adding programme completion
to the adjustment set moves the point estimate from 1420 to 2180 units of annual
earnings, and the second number is further from the experimental benchmark of
Section~\ref{sec:sens-benchmark}, not closer.

\section{Front door adjustment}
\label{sec:graph-frontdoor}

When no admissible backdoor set exists, identification is occasionally still
possible through a fully mediating variable.

\begin{theorem}[Front door adjustment]
\label{thm:frontdoor}
\index{front door criterion}
Let $M$ satisfy: every directed path from $A$ to $Y$ passes through $M$; no
backdoor path from $A$ to $M$ is unblocked; and every backdoor path from $M$ to
$Y$ is blocked by $A$. Then
\begin{equation}
  \E\left[ Y(a) \right] = \sum_m \Prob(M = m \mid A = a)
    \sum_{a'} \E\left[ Y \mid M = m, A = a' \right] \Prob(A = a').
  \label{eq:frontdoor}
\end{equation}
\end{theorem}

\begin{proof}
Apply Theorem~\ref{thm:backdoor} twice. The second condition gives that the
effect of $A$ on $M$ is identified with the empty adjustment set, so
$\Prob(M(a) = m) = \Prob(M = m \mid A = a)$. The third condition gives that the
effect of $M$ on $Y$ is identified adjusting for $A$, so
$\E[Y(m)] = \sum_{a'} \E[Y \mid M = m, A = a'] \Prob(A = a')$. The first
condition gives $Y(a) = Y(M(a))$, and composing the two identified pieces along
the mediator yields~\eqref{eq:frontdoor}.
\end{proof}

The front door criterion is more often cited than used
\citep{weiss2020frontdoor}. Its first condition,
that the mediator captures the entire effect, is as strong as the ignorability
it replaces and is equally untestable. Table~\ref{tab:criteria-summary}
summarises when each strategy applies.

\begin{table}[t]
\centering
\caption[Comparison of identification strategies]{Comparison of the
identification strategies developed in Part~\ref{part:identification} and
Part~\ref{part:design}. The final column records what a violation of the key
assumption does to the estimate, which is the property that matters when the
assumption is doubtful rather than false.}
\label{tab:criteria-summary}
\small
\begin{tabular}{p{26mm}p{28mm}p{13mm}p{32mm}}
\toprule
Strategy & Key assumption & Testable & Failure mode \\
\midrule
Backdoor & no unmeasured confounder & no & bias of unknown sign \\[2pt]
Front door & complete mediation & no & bias toward the null \\[2pt]
Instrument & exclusion restriction & partly & bias amplified by a weak first stage \\[2pt]
Difference in differences & parallel trends & partly & bias equal to the trend gap \\[2pt]
Regression discontinuity & continuity at the cutoff & partly & bias local to the cutoff \\
\bottomrule
\end{tabular}
\end{table}

The column marked partly testable deserves emphasis. An overidentification test\index{overidentification test}
for an instrument, or a pre trend test\index{pre trend test} for difference in differences, has power
against some violations and none against others, and passing it is evidence
rather than proof. Chapter~\ref{chap:sensitivity} takes the position that the
right response to an untestable assumption is not a test but a calibrated
statement of how wrong it would have to be.
